\begin{afterword}
\summaryhead

The author recalls a past self, Eliezer1999, who claimed to know nothing about morality and hoped that a generic superintelligence would work it out. The post argues that this view undercut itself. As Davidson observed, a belief with no correct content is not about anything. A program can compute ``morality'' only if something in its start state points it there, since by the no-free-lunch theorems no procedure is favoured when nothing is known. And a word with nothing built into it, like ``wakalixes,'' means nothing.

A scientist's notion of evidence can change because the brain already contains what new arguments appeal to; in the same way, what your brain would recognize as right is what ``right'' is about. Distrusting every evolved intuition would discard the whole brain, and a morality no human could ever be moved to adopt would be no better than a stone tablet saying ``Thou shalt murder.'' Once some content is built into the question ``What is truly right'', the post asks, why not build in more, such as that joy is preferable to sorrow? Then you know quite a lot about morality, though nothing for certain.

\responsehead

The post's first argument is sound in the form it states. A person who claims to know nothing at all about morality cannot say what the word refers to, so ``I don't know what is right'' becomes as empty as ``I don't know what is wakalixes.'' The post needs only the weak premise that ``at least some'' beliefs must be right before others can be wrong, which is less than Davidson's ``endless true beliefs.'' Its supporting parallels hold up: the no-free-lunch theorems are summarized fairly, and the method it recommends, starting from one's own judgments and revising them, is reflective equilibrium, whose starting judgments Rawls called ``provisional fixed points.'' The beaver example is Rorty's wording, as the notes on ``Truly Part Of You'' record.

The post then goes further, and here it takes a side. What your brain ``would recognize as right'' is what you are talking about, and intuitions such as ``joy is preferable to sorrow'' are part of the question, not only of the answer. This resembles Frank Jackson's moral functionalism, a form of analytic descriptivism, on which moral properties are whatever fits the platitudes of a folk morality refined by reflection.

The main objection to views of this kind is that they turn disagreement into talking past each other. If ``right'' means what my starting points fix, someone who starts elsewhere means something else, and we no longer disagree. Critics of Jackson's view hold that someone who rejects a shared platitude can be making a moral claim, not a conceptual mistake. The post accepts the consequence for a ghost of perfect emptiness, which may be ``asking a different question.'' It does not take up humans whose starting points differ.

The alternative the post rejects is a light ``that in principle you might not be able to perceive at all.'' That is close to Eliezer1999's view, and the post argues fairly against it. It is a stronger view than moral realism as the Stanford Encyclopedia defines it, which requires only that moral claims purport to report facts and that some of them are true.

In short: a sound argument that moral words need some built-in content, with precedents in Davidson and Rawls, extended to a view close to Jackson's moral functionalism, whose main objection, lost disagreement, the post meets for an alien mind and leaves open for humans.
\end{afterword}
